% Companion tables: table_2s.tex (main text), table_4s.tex (appendix).
% Uses the new continuous-load sensitivity matrix only.
\subsection{Parameter Sensitivity}
\label{sec:parameter-sensitivity}

We individually halve or double $n_0$, $\alpha$, $H$, and $\tau$
around their defaults. Table~\ref{tab:parameter-sensitivity-2s}
reports serving capacity and mean TTFT at the same KV budgets and
on the same 24 paired windows per workload. All latency measurements
use a common offered load of 2 QPS. Capacity uses a 2-s mean-TTFT
limit with the same TBT and backlog constraints as the main evaluation.

Threshold strength affects both metrics most clearly under tight KV
budgets. On Conversation with 1,024 blocks, increasing $n_0$ from
8 to 16 reduces capacity from 2.20 to 2.04 QPS and increases mean
TTFT from 1,511 to 1,850 ms. Reducing $\alpha$ to 64 instead yields
2.30 QPS and 1,362 ms, showing that the default is not uniformly
optimal across operating conditions. Temporal parameters have
smaller effects on capacity in this matrix: halving or doubling $H$
changes capacity by $-0.59$\% to $+1.03$\% relative to the default,
and changing $\tau$ yields $-1.78$\% to $+2.73$\%.

The paired metrics also distinguish latency at a fixed load from
the capacity boundary. At 4,096 blocks and 2 QPS, all parameter
configurations yield approximately 734 ms on Code and 332 ms on
Conversation, while their capacity estimates differ. Under a relaxed
4-s limit, for example, Code capacity increases from 3.76 to 4.14 QPS
with $n_0=16$, despite unchanged latency at 2 QPS
(Table~\ref{tab:parameter-sensitivity-4s}). Parameter sensitivity
therefore depends on both available KV space and the service objective.
